\section{BDD Scenarios}
\label{sec:bdd}

Below are BDD scenarios in Gherkin format for key system functions.

\begin{codelisting}
\caption{New User Registration}
\label{lst:bdd}
\begin{skedcode}
Scenario: New user registration
  Given the registration page is open
  And email "test@example.com" is not yet registered in the system
  When the user enters email "test@example.com", a password, and a nickname
    and confirms registration
  Then the user gains access to an empty personal catalog
  And a repeated registration attempt with the same email shows the message
    "An account with this email already exists"
\end{skedcode}
\end{codelisting}
\sked{K7}

\begin{codelisting}
\caption{Account Login}
\label{lst:bdd-login}
\begin{skedcode}
Scenario: Login with correct credentials
  Given the user is registered in the system
  When they enter the correct email and password and click "Sign in"
  Then they are redirected to their personal catalog
  And they see their previously added items
\end{skedcode}
\end{codelisting}
\sked{K8}

\begin{codelisting}
\caption{Logout and Protected Page Access}
\label{lst:bdd-logout}
\begin{skedcode}
Scenario: Accessing a protected page after logout
  Given the user is logged in and viewing the catalog
  When the user clicks "Log out"
  Then they are redirected to the login page
  And navigating directly to /catalog redirects back to the login page
  And logging in again restores full access to the catalog
\end{skedcode}
\end{codelisting}
\sked{K9}

\begin{codelisting}
\caption{Adding a New Item}
\label{lst:bdd-add}
\begin{skedcode}
Scenario: Adding a new content item to the catalog
  Given the user is on the catalog page
  When the user opens the "Add Item" form,
    enters title "Dune" and selects type "Book",
    and submits the form
  Then the item "Dune" of type "Book" appears in the catalog
  And the item displays the default cover illustration for books
  And no other user's catalog is affected
\end{skedcode}
\end{codelisting}
\sked{K13}

\begin{codelisting}
\caption{Deleting an Item with Cascade Removal}
\label{lst:bdd-delete}
\begin{skedcode}
Scenario: Deleting an item removes it from all collections
  Given item "Interstellar" exists in the catalog
  And "Interstellar" belongs to collections "Favourites" and "2025 Films"
  When the user confirms deletion of "Interstellar"
  Then "Interstellar" no longer appears in the catalog
  And "Interstellar" no longer appears in "Favourites" or "2025 Films"
  And both collections still exist and contain their other items
\end{skedcode}
\end{codelisting}
\sked{K15}

\begin{codelisting}
\caption{Repeated Viewing with a Different Rating}
\label{lst:bdd-session}
\begin{skedcode}
Scenario: Re-watching a film with a different rating
  Given content item "Film X" already has one Session with rating 9
  When the user adds a new Session for the same item
    with rating 5 and note "didn't enjoy it the second time"
  Then the content item has two Sessions, each with its own date,
    rating, and note
  And both ratings are stored separately, neither overwrites the other
\end{skedcode}
\end{codelisting}
\sked{K18}

\begin{codelisting}
\caption{Item in Multiple Collections}
\label{lst:bdd-collection}
\begin{skedcode}
Scenario: An item belongs to multiple collections simultaneously
  Given the user has item "Book Y" in the catalog
  And collections "Favourites" and "2026" exist
  When the user adds "Book Y" to both collections
  Then "Book Y" is displayed in the contents of both collections
  And removing "Book Y" from collection "2026" does not remove it
    from "Favourites" or from the catalog
\end{skedcode}
\end{codelisting}
\sked{K21,K25}

\begin{codelisting}
\caption{Automatic Completion Date Fill}
\label{lst:bdd-date}
\begin{skedcode}
Scenario: Automatic completion date fill
  Given the user is editing an item with status "in progress"
  When the user changes the status to "completed"
    and does not enter a date manually
  Then the system automatically sets the completion date
    equal to the current date
\end{skedcode}
\end{codelisting}
\sked{K27}

\begin{codelisting}
\caption{Filtering by Tag}
\label{lst:bdd-tag}
\begin{skedcode}
Scenario: Filtering the catalog by tag
  Given the catalog contains three items: two tagged "fantasy" and one without
  When the user selects the filter by tag "fantasy"
  Then only the two items tagged with that tag are displayed
  And when the filter is cleared, all three items are displayed again
\end{skedcode}
\end{codelisting}
\sked{K17}

\begin{codelisting}
\caption{Searching by Partial Title}
\label{lst:bdd-search}
\begin{skedcode}
Scenario: Searching the catalog by partial title
  Given the catalog contains "Dune", "Dune Messiah", and "Interstellar"
  When the user types "dune" in the search field
  Then only "Dune" and "Dune Messiah" are displayed
  And when the search field is cleared, all three items appear again
\end{skedcode}
\end{codelisting}
\sked{K29}

\begin{codelisting}
\caption{Registration Input Validation}
\label{lst:bdd-validation}
\begin{skedcode}
Scenario: Registration with invalid input is rejected
  Given the registration page is open
  When the user submits the form with password "abc" (fewer than 8 characters)
  Then the form is not submitted
  And the message "Password must be at least 8 characters" is shown
  When the user submits with a malformed email "notanemail"
  Then a message about an invalid email format is shown
  When the user submits with a nickname of 1 character
  Then a message about nickname length (2--50 characters) is shown
\end{skedcode}
\end{codelisting}
\sked{K46}
